\section{Uniform regularity on caps}\label{sec:caps}

The purpose of this section is to turn low-order control of an
admissible height into all-orders control wherever the metric has
all-orders control in a lower or upper cap. The edge facing the
center is artificial: no high-order boundary data will be imposed
there.

Throughout this section we consider families of metrics \(g\in\mathcal U\)
and admissible heights \(z\). Assume fixed uniform bounds
\begin{equation}\label{eq:low-bounds}
\norm g_{C^8(S)}+\norm z_{C^8(S)}\le M,\qquad
g\ge cI,\qquad \abs{H_{yy}}\ge c,\qquad E\ge c,
\end{equation}
where \(M<\infty\) and \(c>0\). Constants in this section may depend
on these bounds, on \(\kappa,h\), and on the derivative order.

Set
\begin{equation}\label{eq:DGa}
D=\pd_x-q\pd_y,\qquad G=\frac E{H_{yy}^2},\qquad a=KG.
\end{equation}
Use flow coordinates \(t,s\) defined by
\begin{equation}\label{eq:flow}
x=t,\qquad y_t(t,s)=-q(t,y(t,s)),\qquad y(0,s)=s.
\end{equation}
They exist on \([-2,2]^2\), have image in the interior of \(S\),
and satisfy
\begin{equation}\label{eq:flow-bounds}
\abs{y(t,s)-s}\le1/50,\qquad D=\pd_t,\qquad
0<c_0\le s_y\le C_0.
\end{equation}
Indeed the first estimate follows from \(\abs q\le1/100\).
Differentiating the flow gives
\[
y_s=\exp\left(-\int_0^t q_y(r,y(r,s))\dd r\right)>0.
\]
The low bounds give uniform bounds for the derivatives of the chart
and its inverse through the orders used below, in particular through
order three. Denote this chart by \(\Phi_z(t,s)=(t,y(t,s))\).

\begin{proposition}[Cap regularity]\label{prop:cap}
Fix \(-1/4<b_*\le0\), and assume \eqref{eq:low-bounds}.
Suppose that, for every compact rectangle
\[
R\Subset(-2,2)\times(-2,b_*),
\]
the ordinary coordinate derivatives of \(g\) of every order have
uniform bounds on \(\Phi_z(R)\).
Then the ordinary coordinate derivatives of \(z\) of every order
also have uniform bounds on \(\Phi_z(R)\).
The corresponding assertion holds for the upper cap
\((-2,2)\times(-b_*,2)\).
\end{proposition}

Here and below a rectangle has sides parallel to the \(t,s\) axes.
Its image may depend on the solution, but all margins in coordinate
labels are chosen independently of the member of the family.
More precisely, if \(R\Subset R'\) are two such rectangles, the
uniform bounds for the chart and its inverse give a fixed radius
of ordinary coordinate balls around points of \(\Phi_z(R)\)
that stay inside \(\Phi_z(R')\). We can therefore apply local
Sobolev estimates with uniform constants in the original coordinates.
No high derivatives of the chart are needed for this step.

\subsection{The differentiated equation}

At the induction step we estimate \(u=\pd_y^lz\) first; the
Darboux equation then recovers derivatives with more \(x\) indices.
We differentiate in the original coordinates before passing to
the flow chart, whose high derivatives are not yet controlled.
Solving the Darboux equation for \(z_{xx}\) gives
\begin{equation}\label{eq:P}
z_{xx}=P(x,y,z_i,z_{xy},z_{yy})
=\Gamma^r_{xx}z_r+\frac{H_{xy}^2+KE}{H_{yy}}.
\end{equation}
The metric and its derivatives are prescribed arguments of \(P\).
In particular,
\begin{equation}\label{eq:P-partials}
P_{z_{xy}}=2q,\qquad P_{z_{yy}}=-q^2-a,\qquad
P_{z_i}=\Gamma^i_{DD}+O(K).
\end{equation}
The notation \(\Gamma^i_{DD}\) contracts the two lower Christoffel
indices against \(D\), without differentiating \(D\).
The last identity follows more explicitly from
\[
P_{z_i}=\Gamma^i_{DD}
+K\left(\frac{E_{z_i}}{H_{yy}}
+\frac{E\Gamma^i_{yy}}{H_{yy}^2}\right).
\]
All coefficients in an \(O(K)\) term in this section are uniformly
bounded under \eqref{eq:low-bounds}.

\begin{lemma}[Curvature drift]\label{lem:drift}
The quantities in \eqref{eq:DGa} satisfy
\begin{equation}\label{eq:drift}
Dq-P_{z_x}q-P_{z_y}=K_yG+O(K).
\end{equation}
\end{lemma}

\begin{proof}
Put \(Y=\pd_y\). The definitions and the Darboux equation give
\[
H(Y,D)=0,\qquad H(D,D)=KE/H_{yy},
\qquad H(D,\pd_x)=KE/H_{yy}.
\]
Thus \(H(D,\cdot)=O(K)\) on bounded coordinate vectors.
The commutation formula for a covariant Hessian gives
\begin{equation}\label{eq:hessian-commutation}
(\nabla_DH)(Y,D)=(\nabla_YH)(D,D)+O(K),
\end{equation}
because in dimension two the curvature tensor is a bounded tensor
multiple of \(K\), and the commutator contracts it against \(dz\).

Differentiate \(H(Y,D)=0\). The term containing \(\nabla_DY\)
is \(O(K)\), and
\(\nabla_DD=\Gamma^i_{DD}\pd_i-(Dq)Y\). Hence the left side
of \eqref{eq:hessian-commutation} equals
\[
H_{yy}(Dq-q\Gamma^x_{DD}-\Gamma^y_{DD})+O(K).
\]
On the right, differentiating \(H(D,D)=KE/H_{yy}\) gives
\(K_yE/H_{yy}+O(K)\); the two connection terms are \(O(K)\).
Division by \(H_{yy}\), followed by \eqref{eq:P-partials},
proves \eqref{eq:drift}.
\end{proof}

\begin{lemma}[High derivative equation]\label{lem:high-equation}
Let \(l\ge8\). Suppose that the ordinary derivatives of \(z\) through
order \(l\) have uniform \(L^2\) bounds on images of compact rectangles
in the lower cap, and that the metric satisfies the hypothesis of
\Cref{prop:cap}. Then \(u=\pd_y^l z\), re-expressed in flow coordinates,
satisfies
\begin{equation}\label{eq:linear-cap}
u_{tt}+Au_{ss}+Bu_t+Cu_s=R_l,
\end{equation}
where \(R_l\) is uniformly bounded in \(L^2\) on smaller compact
rectangles, and
\begin{equation}\label{eq:linear-coefficients}
\begin{gathered}
A=G_1K,\qquad G_1=Gs_y^2,\qquad
B=-P_{z_x}-2lq_y,\\
C=(l+1)A_s+O_l(A).
\end{gathered}
\end{equation}
The function \(G_1\) has positive lower and upper bounds and bounded
first derivatives. The coefficients \(A,A_t,A_s,B,B_s,C\) are
uniformly bounded.
\end{lemma}

\begin{proof}
The inductive \(L^2\) bounds give supremum bounds through order
\(l-2\) on smaller sets by the local two-dimensional Sobolev
inequality on the ordinary coordinate balls described above.

Differentiate \eqref{eq:P} \(l\) times in the original \(y\)
coordinate. In addition to the direct linearization, the terms
containing derivatives of \(z\) of order \(l+1\) are
\[
l\bigl(2q_yu_x-(q^2+a)_yu_y\bigr).
\]
These arise when a second-jet argument receives \(l-1\)
differentiations and the remaining single derivative is placed
elsewhere. The coefficient derivatives are total derivatives,
including the explicit metric dependence of \(P\).

Every remaining differentiated solution factor has order at most
\(l\). A factor of order at least \(l-1\) consumes at least \(l-3\)
of the \(l\) differentiations. There cannot be two such factors,
since \(2(l-3)>l\) for \(l\ge8\). Thus all but at most one factor
use the supremum bounds through order \(l-2\); the remaining factor
uses its \(L^2\) bound. Metric derivatives are bounded by hypothesis,
and inverse powers of \(H_{yy}\) are bounded by
\eqref{eq:low-bounds}. The remainder is consequently bounded in
\(L^2\).

Using
\[
D^2=\pd_x^2-2q\pd_{xy}+q^2\pd_y^2-(Dq)\pd_y
\]
and \(\pd_x=D+q\pd_y\), the differentiated equation becomes
\begin{equation}\label{eq:pre-flow}
\left[
D^2+a\pd_y^2+BD+
\bigl(Dq-P_{z_x}q-P_{z_y}+la_y\bigr)\pd_y
\right]u=R_l.
\end{equation}
Now \(D=\pd_t\), \(\pd_y=s_y\pd_s\), and
\(\pd_y^2=s_y^2\pd_s^2+s_{yy}\pd_s\).
By \Cref{lem:drift}, the transverse coefficient is
\[
C=a s_{yy}+s_y(K_yG+la_y+O(K))
=(l+1)Gs_yK_y+O_l(K).
\]
Since
\[
A_s=Gs_yK_y+K(Gs_y^2)_s,
\]
this is \eqref{eq:linear-coefficients}. The factor \(l+1\) has two
sources: the commutation identity contributes one copy of
\(Gs_yK_y\), and the term \(la_y\) contributes \(l\) copies.
Every other transverse term contains \(K\) times a bounded factor.
All asserted low coefficient
bounds follow from \eqref{eq:low-bounds} and the flow bounds.
\end{proof}

\subsection{A weighted gradient estimate}

The differentiated equation reduces the cap induction to a gain
of one derivative. Its principal part changes type, so an elliptic
estimate alone is insufficient. We combine a directed multiplier
near and inside the negative region with a weighted elliptic
estimate on its complement.

\begin{lemma}[The free cap edge]\label{lem:gradient}
Retain the model geometry and flow coordinates fixed above.
Let \(l\ge8\), and let the coefficients satisfy
\eqref{eq:linear-coefficients} with the uniform bounds in
\Cref{lem:high-equation}. For every
\(R_0\Subset(-2,2)\times(-2,b_*)\), there is a larger compact
rectangle \(R_1\) in that cap and a constant \(C_{R_0,l}\) such that
every smooth solution of
\[
u_{tt}+Au_{ss}+Bu_t+Cu_s=f
\]
on a neighborhood of \(R_1\) satisfies
\begin{equation}\label{eq:gradient-bound}
\norm{(u_t,u_s)}_{L^2(R_0)}
\le C_{R_0,l}\bigl(\norm u_{L^2(R_1)}+\norm f_{L^2(R_1)}\bigr).
\end{equation}
The constants are uniform over the families under consideration.
\end{lemma}

\begin{proof}
Choose \(b<b_*\) above the upper edge of \(R_0\).
Integrations will take place below \(s=b\), with their other three
edges sufficiently far out that cutoff derivatives lie in
\[
3/2<\abs t<2
\quad\text{or}\quad -2<s<-3/2.
\]
Both regions have uniformly positive curvature. At the lateral
edges, \(t^2-h(y)\ge t^2-1>0\). At the lower edge,
\eqref{eq:flow-bounds} gives \(y<-1\), so \(h(y)<0\).
There is room for nested cutoffs there and for a compact rectangle
\(R_1\) containing every integration set, including its upper edge.

Set \(d=b-s\) and \(p=8\). We choose the directed weight to vanish
like \(d^p\) at the artificial edge; it remains positive on \(R_0\).
Differentiating such a weight costs a factor \(d^{-1}\), so the
two estimates will meet at \(K=c_1d\), for a small constant
\(c_1>0\) chosen below. The directed multiplier controls
\(K\le c_1d\), where its positive-curvature error is
\(O(K/d)\). On \(K\ge c_1d\), an elliptic estimate controls
the stronger weight \(d^{p-1}\), which pays for the weight and
cutoff derivatives. \Cref{fig:cap} shows the integration geometry.

\begin{figure}[H]
\centering
\begin{tikzpicture}[font=\small,>=Latex,xscale=1.35,yscale=1.2]
  \fill[blue!6] (-2,-2) rectangle (2,0);
  \fill[orange!22]
    (-1.12,0)--(-1.12,-.35)
    .. controls (-1.12,-.83) and (-.57,-1.28) .. (0,-1.28)
    .. controls (.57,-1.28) and (1.12,-.83) .. (1.12,-.35)
    --(1.12,0)--cycle;
  \draw[thick]
    (-1.12,0)--(-1.12,-.35)
    .. controls (-1.12,-.83) and (-.57,-1.28) .. (0,-1.28)
    .. controls (.57,-1.28) and (1.12,-.83) .. (1.12,-.35)
    --(1.12,0);
  \draw[very thick,blue!65!black] (-2,0)--(-2,-2)--(2,-2)--(2,0);
  \draw[thick,dashed] (-2,0)--(2,0);
  \node[above=7pt] at (0,0) {artificial edge \(s=b\)};
  \node at (-.2,-.48) {\(K<0\)};
  \node at (0,-1.65) {\(K>0\)};
  \node[anchor=east] at (-1.15,-.85) {\(K=0\)};
  \draw[<->] (.66,-.04)--(.66,-.82);
  \node[anchor=west] at (.72,-.58) {\(d\)};
  \draw[->,gray] (2.45,-1.8)--(2.45,-.65) node[above,black] {\(s\)};
  \draw[->,gray] (2.2,-1.55)--(3.05,-1.55) node[right,black] {\(t\)};
  \node[text=blue!65!black] at (0,-2.5) {three edges in positive curvature};
\end{tikzpicture}
\caption{A lower-cap integration region in flow coordinates, schematically
and not to scale. The lateral and bottom edges are chosen in positive
curvature. The artificial edge \(s=b\) may cross the negative region;
no high-order data are prescribed there. The distance \(d=b-s\)
enters the vanishing weight. The actual curvature zero set depends on
the solution through the flow coordinates; only its position relative
to the three elliptic edges and the artificial edge is illustrated.}
\label{fig:cap}
\end{figure}


\smallskip
\noindent\emph{The directed multiplier.}
The transverse multiplier uses the favorable curvature drift.
We add a small \(t\)-directed part to obtain positivity at those
zeros of curvature where \(K_y=0\): there the \(x^2\) term in
the model curvature must supply the estimate.
Put
\begin{equation}\label{eq:weights}
\begin{gathered}
I(t,s)=\int_0^tB(r,s)\dd r,\\
W=d^p e^{-\lambda s+I(t,s)},\qquad
m=-W,\qquad n=\epsilon tW,\qquad H_0=\lambda+p/d .
\end{gathered}
\end{equation}
Here \(\epsilon>0\) will be small and \(\lambda>0\) large.
The factor \(e^I\) gives \(m_t=Bm\), canceling the corresponding
mixed term in the identity below. Also
\(m_s/W=H_0-I_s\), with \(I_s\) uniformly bounded.

Testing against \(mu_s+nu_t\) and integrating by parts gives
the interior quadratic form
\begin{equation}\label{eq:quadratic}
T u_t^2+S_0u_s^2+Ju_tu_s,
\qquad
\begin{cases}
T=(m_s-n_t)/2+Bn,\\
S_0=-(Am)_s/2+(An)_t/2+Cm,\\
J=-m_t+Bm-(An)_s+Cn.
\end{cases}
\end{equation}
The symbol \(S_0\) denotes this coefficient, not the model square.
Substituting \eqref{eq:weights} gives, for large \(\lambda\),
\begin{equation}\label{eq:quadratic-bounds}
\begin{aligned}
T/W&\ge H_0/4,\\
S_0/W&=-(l+1/2)A_s+\epsilon tA_t/2-H_0A/2+O_l(A),\\
\abs J/W&\le C_l\epsilon(1+H_0\abs A).
\end{aligned}
\end{equation}

\smallskip
\noindent\emph{Positivity near and inside the negative region.}
There are choices of these constants and \(c_1,c_2>0\) such that
\begin{equation}\label{eq:coercive}
T u_t^2+S_0u_s^2+Ju_tu_s
\ge c_2W(u_t^2+u_s^2)
\qquad\text{where }K\le c_1d.
\end{equation}
We verify the choices in detail.

For \(\abs K\) sufficiently small, \eqref{eq:negative-box}
excludes \(\mathcal P\), so the prescribed curvature formula applies.
In the lower cap \(y\le1/50\). The plateau and monotonicity in
\eqref{eq:h} imply \(K_y\le0\) there. Also
\begin{equation}\label{eq:A-derivatives}
A_s=G_1K_y/s_y+O(A),\qquad
A_t=G_1(2\kappa t-qK_y)+O(A).
\end{equation}
Thus the first two terms in the second line of
\eqref{eq:quadratic-bounds}, up to \(O_l(A)\), equal
\[
G_1\left[
-\left(\frac{l+1/2}{s_y}+\frac{\epsilon tq}{2}\right)K_y
+\epsilon\kappa t^2
\right].
\]
For small fixed \(\epsilon>0\) this is bounded below by a positive
constant times \(-K_y+\epsilon t^2\). To make its uniform
positivity independent of the flow chart, consider the fixed
physical zero set
\[
\mathcal Z=\{(x,y)\in[-2,2]\times[-101/50,1/50]:x^2=h(y)\}.
\]
It is compact and contains every curvature zero in the lower cap.
If \(x=0\) on \(\mathcal Z\), then \(y=-1\), where \(K_y<0\);
at every other point \(x^2>0\). Thus
\(-K_y+\epsilon x^2\) has a positive minimum on \(\mathcal Z\).
The uniform bounds for \(G_1,s_y,q\) transfer this minimum to
the geometric contribution on a common neighborhood of \(K=0\).

The cross-term loss after absorbing part of \(T u_t^2\) is bounded by
\begin{equation}\label{eq:cross-loss}
\frac{J^2}{TW}
\le C_l'\epsilon^2\left(H_0^{-1}+H_0A^2\right).
\end{equation}
On \(K\le0\), the term \(-H_0A/2=H_0\abs A/2\) is favorable.
Since \(\abs A\) is bounded, choose \(\epsilon\) small enough
to absorb the second term of \eqref{eq:cross-loss} into a small
fraction of \(H_0\abs A\). Increase \(\lambda\) to absorb
the \(O_l(A)\) terms and the \(H_0^{-1}\) loss.
Away from a neighborhood of the zero set on the negative side,
the same increase of \(\lambda\) makes \(H_0\abs A\) dominate
all bounded adverse terms.

Fix these \(\epsilon,\lambda\). On \(0<K\le c_1d\), since \(d\)
is bounded above,
\[
H_0A\le Cc_1(\lambda d_{\max}+p).
\]
Choose \(c_1\) so small that this region lies in the preceding
neighborhood of the zero set, and that the adverse \(H_0A\)
term and remaining errors use only a fraction of its positive
lower bound. This proves \eqref{eq:coercive}. The order of
choices is \(\epsilon\), then \(\lambda\), then \(c_1\).

\smallskip
\noindent\emph{The auxiliary elliptic estimate.}
On the remaining region we claim
\begin{equation}\label{eq:aux-elliptic}
\iint_{\{K\ge c_1d\}} d^{p-1}(u_t^2+u_s^2)\dd t\dd s
\le C\bigl(\norm u_{L^2(R_1)}^2+\norm f_{L^2(R_1)}^2\bigr),
\end{equation}
on any of the supports needed for the directed multiplier.
To prove it, choose a cutoff \(\chi\) farther out at the three
elliptic edges, equal to one on those supports. Let \(\phi\) be
a smooth cutoff equal to zero on \((-\infty,c_1/2]\) and one
on \([c_1,\infty)\), and test the equation against \(-vu\), where
\[
v=d^{p-2}\chi^2\phi(K/d)^2.
\]
On \(v>0\), \(A\ge c'd\). In the transition region of \(\phi\),
\(K=O(d)\), and hence derivatives of \(K/d\) are \(O(d^{-1})\).
It follows that
\[
\abs{v_t}+\abs{v_s}
\le C\sqrt v\,d^{(p-2)/2}(1+d^{-1}).
\]
The squared cutoffs make this inequality valid also at their zeros.

Integration by parts produces the positive term
\(\iint v(u_t^2+Au_s^2)\) and the error
\[
\iint u\bigl[(v_t-vB)u_t+((Av)_s-vC)u_s\bigr].
\]
On \(v>0\), the coefficient bounds and \(A\ge c'd\) give
\[
\frac{(v_t-vB)^2}{v}
+\frac{((Av)_s-vC)^2}{Av}
\le C(1+d^{p-5}).
\]
Young's inequality absorbs half of the positive term; the remaining
cost is at most \(C\iint(1+d^{p-5})u^2\).
Since \(p=8\), this is bounded by \(C\norm u_2^2\).
The source term is at most \(C\norm f_2\norm u_2\).
We conclude that
\[
\iint v(u_t^2+Au_s^2)
\le C(\norm u_2^2+\norm f_2^2).
\]
Where \(\chi=\phi=1\), \(v=d^{p-2}\gtrsim d^{p-1}\)
and \(vA\gtrsim d^{p-1}\), proving \eqref{eq:aux-elliptic}.
This estimate permits ellipticity to degenerate linearly as
\(s\) approaches \(b\).

\smallskip
\noindent\emph{Combining the estimates.}
Multiply \(m,n\) by one nonnegative cutoff toward the other three
edges, equal to one on \(R_0\) and on the nonelliptic portion of
the integration region. Its derivative terms lie in fixed positive
curvature; decreasing \(c_1\), if needed, puts those supports
in the region controlled by \eqref{eq:aux-elliptic}.
The coefficients in \eqref{eq:quadratic} and the cutoff errors
are \(O(d^{p-1})\), since \(W\asymp d^p\) and \(H_0=O(1+d^{-1})\).
Thus \eqref{eq:aux-elliptic} controls all adverse terms outside
\eqref{eq:coercive}.

On the coercive region the source \(f(mu_s+nu_t)\) is absorbed by
Young's inequality, at a cost \(C\norm f_2^2\). On the complementary
region use \eqref{eq:aux-elliptic} and
\(W^2/d^{p-1}=O(d^{p+1})\) to obtain the same bound together
with \(C\norm u_2^2\). Combining the coercive part with the
elliptic estimate on its complement, we obtain
\[
\iint_{R_0}W(u_t^2+u_s^2)
\le C(\norm u_{L^2(R_1)}^2+\norm f_{L^2(R_1)}^2).
\]
The weight is bounded below on \(R_0\), so this proves
\eqref{eq:gradient-bound}.

All integrations at the artificial edge can first be made below
\(s=b-\eta\) and then continued to \(\eta=0\).
For each individual smooth solution, the boundary terms vanish
because \(W\) and \(v\) vanish to positive order. The edge \(b\)
is strictly inside the smooth domain. No uniform bound on its
high derivatives is needed to obtain these identities; the uniform
estimates involve only the interior \(L^2\) norms just displayed.
\end{proof}

\subsection{Completion of the cap induction}

\begin{proof}[Proof of \Cref{prop:cap}]
We prove uniform \(L^2\) bounds for ordinary derivatives of \(z\)
through order \(l\), on images of all compact rectangles in the
cap. The case \(l=8\) follows from \eqref{eq:low-bounds}.
Assume the assertion at some \(l\ge8\). Apply
\Cref{lem:high-equation,lem:gradient} with \(f=R_l\).
The required bounds for \(u,R_l\) on the larger compact rectangle
are supplied by the induction hypothesis. Measures in the two
coordinate systems are uniformly comparable.

The resulting gradient bound controls
\[
\pd_y^{l+1}z=s_yu_s,\qquad
\pd_x\pd_y^lz=u_t+q s_yu_s.
\]
To recover \(\pd_x^k\pd_y^{l+1-k}z\) for \(k\ge2\), apply
\(\pd_x^{k-2}\pd_y^{l+1-k}\) to \eqref{eq:P}.
The two top terms on the right are
\[
P_{z_{xy}}\pd_x^{k-1}\pd_y^{l+2-k}z
\quad\text{and}\quad
P_{z_{yy}}\pd_x^{k-2}\pd_y^{l+3-k}z.
\]
They have fewer \(x\) indices. All remaining solution factors
have order at most \(l\) and obey the same Sobolev product bounds
as in \Cref{lem:high-equation}. Increasing \(k\) therefore gives
every derivative of order \(l+1\).

This closes the induction. For each fixed derivative order, local
Sobolev estimates on slightly larger compact rectangles convert
the \(L^2\) bounds into the asserted supremum bounds. The final
rectangle is arbitrary and fixed; only the constants and intermediate
margins depend on the derivative order.
Reflection \(y\mapsto-y\), using the evenness of \(h\), proves
the upper-cap assertion.
\end{proof}
